\chapter{HTTP/1.1 on the wire}
\chapterintro{HTTP semantics are abstract; HTTP/1.1 gives them a textual start line, fields, framing rules, and connection behavior. Small parsing shortcuts here can become security vulnerabilities.}

\begin{objectives}
\begin{itemize}
  \item Read request lines, status lines, header sections, and message bodies as bytes.
  \item Determine message-body length for the subset implemented in this book.
  \item Explain persistent connections and why exact consumption matters.
  \item Recognize ambiguous framing and request-smuggling risk.
\end{itemize}
\end{objectives}

\section{Message shape}

An HTTP/1.1 message contains a start line terminated by CRLF, zero or more field
lines each terminated by CRLF, an empty CRLF line, and an optional body. A
request uses a request line; a response uses a status line.

\begin{wiremessage}
POST /tasks?notify=false HTTP/1.1
Host: 127.0.0.1:8080
Content-Type: application/json
Content-Length: 21

{"title":"Read RFC"}
\end{wiremessage}

The displayed blank line is mandatory. In actual bytes, every line ends with
\code{\\r\\n}. The body is 21 UTF-8 bytes in this example. A parser must not
search for a newline to end JSON; JSON can contain escaped newlines and the next
request may follow immediately.

\section{The request line}

The request line contains method, one space, request target, one space, and HTTP
version. For a direct request to an origin server, the common target form is the
path and optional query, called \term{origin-form}. Proxies can receive an
absolute URI. CONNECT and OPTIONS define other forms.

Our server accepts HTTP/1.1 origin-form targets and the asterisk target for
server-wide OPTIONS. It rejects other forms explicitly. It also preserves method
case so a lowercase \code{get} is not silently turned into the standard GET.

\section{Field parsing}

Field names use a restricted token grammar and are case-insensitive. A colon
separates name from value. Obsolete line folding is rejected. Control characters
in values are rejected. Exactly one non-empty \code{Host} field is required for
the implemented HTTP/1.1 subset.

Repeated fields need field-specific handling. The final \code{Headers} type
stores an ordered list and exposes both \code{get} and \code{get\_all}. For
framing, the implementation deliberately rejects repeated \code{Content-Length}
even when values match. This is stricter than some permitted normalization and
reduces ambiguity in an educational parser.

\section{Where does the body end?}

Message framing is not the same as media parsing. The server must first isolate
the exact body bytes. Only then should JSON code interpret them.

For requests in our subset:

\begin{enumerate}
  \item A \code{Transfer-Encoding} field is rejected because transfer codings,
        including chunked coding, are not implemented.
  \item A simultaneous \code{Transfer-Encoding} and \code{Content-Length} is
        rejected as ambiguous.
  \item More than one \code{Content-Length} is rejected.
  \item One valid non-negative \code{Content-Length} frames exactly that many bytes.
  \item Without either field, the request has a zero-length body. Endpoints that
        require content can answer \status{411 Length Required}.
\end{enumerate}

The parser keeps bytes after the framed body. On a persistent connection those
bytes can begin the next request. Discarding them would corrupt the stream;
including them in the first body would do the same.

\begin{securityboundary}
Request smuggling occurs when two HTTP recipients disagree about where one
request ends and another begins. Conflicting length fields, transfer-coding
differences, and permissive whitespace parsing are common ingredients. A small
server behind a proxy must not assume ambiguity is harmless merely because its
own parser chooses one interpretation.
\end{securityboundary}

\section{Chunked transfer coding}

Chunked coding represents a body as hexadecimal-size chunks followed by a
zero-size chunk and optional trailer fields. It lets a sender begin before the
final content length is known. A compliant general-purpose HTTP/1.1 parser needs
careful support for its grammar, extensions, decoded-size limits, trailers, and
interactions with other fields.

The book explains chunking but the framework returns \status{501 Not Implemented}
and closes the connection whenever \code{Transfer-Encoding} is present. This is
an explicit product boundary, not an accidental omission hidden behind incorrect
parsing.

\section{Response framing and special cases}

A response status line contains version, three-digit status, and an optional
reason phrase. Reason phrases are for humans and do not define status semantics.
Our serializer adds \code{Date}, \code{Server}, \code{Connection}, and usually
\code{Content-Length}.

Responses to HEAD carry the fields that a corresponding GET would produce but
do not carry the content bytes. The \code{Content-Length} can therefore describe
the hypothetical GET body while no body follows. Informational responses,
\status{204 No Content}, and \status{304 Not Modified} have rules that prohibit
message content; our serializer suppresses it.

\section{Persistent connections}

HTTP/1.1 uses persistent connections by default. After a framed response, the
connection can carry another request. Either side can signal
\code{Connection: close}. Persistence avoids repeated TCP and TLS handshakes,
but it also means a worker can be occupied by an idle peer.

The final server supports sequential requests with a five-second idle timeout
and at most 100 requests per connection. It always sends a length-delimited
response when content is allowed. Early stages send \code{Connection: close} so
only one message needs to be considered while fundamentals are introduced.

HTTP/1.1 pipelining permits a client to send later requests without waiting for
earlier responses, with responses kept in order. It proved difficult to use
robustly in browsers and is not a focus here. Our incremental buffer can retain
an already-arrived next request, but each connection processes requests
sequentially.

\section{Limits are protocol behavior}

Without explicit limits, a peer can send an enormous target, unending fields, or
a declared body too large to store. The final configuration limits the request
line to 8 KiB, the header section to 32 KiB and 100 fields, and the body to 1 MiB.
It answers with \status{414 URI Too Long}, \status{431 Request Header Fields Too
Large}, or \status{413 Content Too Large} as appropriate.

\begin{pitfall}
Calling a single receive operation with a 32 KiB buffer does not impose a valid
32 KiB header limit. The call can return fewer bytes, and a peer can continue.
The implementation must accumulate and count until it finds the delimiter or
crosses the limit.
\end{pitfall}

\summarytitle

HTTP/1.1 uses CRLF-delimited start and field lines, a blank line, and explicit
body-framing rules. Persistent connections make exact byte consumption mandatory.
A strict subset can be correct when it clearly rejects unsupported framing and
ambiguous inputs; a permissive half-implementation is dangerous.

\reviewtitle
\begin{enumerate}
  \item Why does JSON parsing begin only after HTTP body framing succeeds?
  \item What happens to bytes after the declared request body?
  \item Why does a HEAD response sometimes contain a nonzero Content-Length?
  \item Why does this server reject transfer coding instead of ignoring it?
\end{enumerate}

\begin{exercise}[Frame two requests]
Write two GET requests consecutively with no extra separator beyond the first
request's final CRLF pair. Mark the exact boundary. Repeat with a POST whose body
ends immediately before the second request.
\end{exercise}

\begin{exercise}[Ambiguity review]
Analyze a request containing both \code{Transfer-Encoding: chunked} and
\code{Content-Length: 4}. Explain how disagreement between a proxy and origin
could cause the following bytes to be interpreted differently.
\end{exercise}

